\ \\[-10cm]
La part anterior d'aquesta guia d'estudi ha explorat alguns dels principals reptes de comunicar múltiples ordinadors.
\concept{Internet} és una solució que funciona per a aquests reptes: la que els humans han creat durant l'últim mig segle,
ara central en moltes activitats socials, comercials i polítiques. 
% 
Aquesta part estudia les principals funcionalitats que ofereix Internet, així com el \conceptRef{suit}{conjunt} de \conceptRef{protocol}{protocols} 
emprats per aconseguir-les.


\chapterimage{earth.jpg}

\chapter{Internet}\label{sec:internet}

Les comunicacions modernes es basen en \conceptRef{stack}{piles} de xarxa organitzades en capes. 
Potser acabes de descarregar aquesta guia d'Internet.
% 
Quines capes s'utilitzen exactament? Quins protocols s'empren, i amb quina finalitat?
% 
Aquest capítol ofereix una introducció per respondre aquestes preguntes.\\[-1cm]

\section{La pila TCP/IP}


\concept{Internet} és una solució pràctica al problema de la comunicació global entre ordinadors,
el desenvolupament de la qual va començar a finals dels anys seixanta i principis dels setanta, i encara continua.

La commutació de \conceptRef{packet}{paquets} es va concebre a principis dels anys seixanta, i la primera 
implementació d'àrea extensa, \concept{ARPANET}, va començar el 1969. 
% 
Un repte clau en aquell moment era interconnectar les \conceptRef{network}{xarxes} existents, 
propietat de diferents organitzacions i basades en diferents tecnologies i \conceptRef{protocol}{protocols}.
% 
Perquè tot fos viable, calia complir tots els requisits següents:
\begin{itemize}
    \item Les xarxes existents havien de funcionar sense actualitzacions de la infraestructura física 
      ni dels seus protocols, només de la manera com s'utilitzen aquests protocols.
    \item Les xarxes existents havien de continuar gestionant-se de manera independent a nivell local (organització interna autònoma).
    \item Els paquets no sempre s'entreguen: alguns es perden.
    \item Els missatges que arriben poden no estar en ordre, i poden estar repetits.
\end{itemize}

La solució adoptada per ARPANET --i més tard per tot el món-- 
és el conjunt de protocols \concept{TCP/IP}, descrit per primera vegada el 1974. 
És un enfocament per capes, dissenyat per ``endollar-se'' damunt de 
les \conceptRef{LAN}{LANs} (xarxes d'àrea local) existents, amb l'objectiu d'oferir capacitats 
de xarxa a totes les aplicacions d'usuari.

\begin{center}
\begin{bytefield}{16}
\begin{rightwordgroup}[curlyshrinkage=10pt]{User programs}
\wordbox{3}{Layer 7: Application}
\end{rightwordgroup} \\
\begin{rightwordgroup}[curlystyle=\color{color1},curlyshrinkage=10pt]{{\color{color1} TCP/IP}}
\wordbox{1}{Layer 4: Transport}\\
\wordbox{1}{Layer 3: Internet}
\end{rightwordgroup} \\
\begin{rightwordgroup}[curlyshrinkage=10pt]{Existing/future LAN technologies}
\wordbox{1}{Layer 2: Network}\\
\wordbox{1}{Layer 1: Physical}
\end{rightwordgroup}
\end{bytefield}
\end{center} 
 
En poques paraules, els principals serveis que ofereixen aquestes capes són els següents:
\begin{enumerate}
\item[\textbf{Layer 1}] \textbf{(Física)}: 
  Enviar dades \concept{digital}s a través de mitjans \conceptRef{analog}{analògics}.
\item[\textbf{Layer 2}] \textbf{(Xarxa)}:
  Identificar i comunicar dispositius dins de la mateixa \concept{LAN}.\\
  De vegades s'utilitza el terme capa d'``enllaç de dades'' (del model OSI).
\item[\textbf{Layer 3}] \textbf{(Internet)}:
  Identificar i comunicar dispositius a través de totes les \conceptRef{LAN}{LANs}.\\
  De vegades s'utilitza el terme capa de ``xarxa'' (també del model OSI).
\item[\textbf{Layer 4}] \textbf{(Transport)}:
  Identificar i comunicar programes (serveis).\\
  A més (tria'n \textit{només} una):
  \begin{enumerate}[label=\alph*)]
   \item Assegurar que un \conceptRef{stream}{flux} de dades arriba correctament a l'altre extrem, 
   i en ordre \\-- \textit{però} s'introdueix una \conceptRef{overhead}{sobrecàrrega} 
   i cal establir \conceptRef{connection}{connexions} (\concept{TCP}).
     
   \item Utilitzar missatges simples i relativament eficients, amb una sobrecàrrega petita\\
   -- \textit{però} els missatges es poden perdre, desordenar o duplicar (\concept{UDP}).
  \end{enumerate}
\item[\textbf{Layer 7}] \textbf{(Aplicació)}:
  Fer qualsevol cosa que l'usuari pugui necessitar. \\
  No t'oblidis de la seguretat (\eg, criptografia) 
  ni de l'eficiència (\eg, compressió): la capa 4 i les inferiors no ho proporcionen.
\end{enumerate}

\section{Visió general dels protocols TCP/IP}

Es defineixen múltiples \conceptRef{protocol}{protocols} per implementar la funcionalitat
de cada \conceptRef{layer}{capa}. Aquestes definicions inclouen el \conceptRef{packet format}{format de paquet} requerit,
les regles per intercanviar aquests paquets i la seva semàntica.
% 
Aquests protocols s'han d'entendre com un sistema, en el qual cada part té tasques molt específiques.

Alguns dels protocols més rellevants per a les comunicacions a través d'Internet 
es mostren a la figura següent. En aquesta:
\begin{itemize}
\item Les fletxes contínues indiquen que un protocol s'encapsula dins d'un altre. 
  Si el destí és una capa, cal utilitzar un dels protocols d'aquesta capa.

\item Les caixes discontínues indiquen protocols auxiliars definits per ajudar altres funcions. 
Les fletxes discontínues indiquen quina capa utilitza quin protocol auxiliar.

\item Les caixes amb línies més gruixudes indiquen els protocols destacats en aquesta guia d'estudi.
\end{itemize}

\includegraphics[width=\linewidth]{protocol_dependencies.pdf}

\begin{exercise}
Volem enviar la frase ``Perdó, arribaré 5 minuts tard'' per correu electrònic.
El nostre client de correu enviarà aquestes dades amb el protocol SMTP, que al seu torn
utilitza TCP (part de la capa 4).
% 
Suposant que SMTP utilitza capçaleres a la seva \concept{PDU}, quantes capçaleres 
contindrà la \concept{PDU} corresponent de la capa 2?
\end{exercise}

\begin{exercise}
Per què creus que la capa~3 només té un protocol (dos si comptes IPv4 i IPv6 per separat)? 
\end{exercise}

\section{Qui mana a Internet?}

La \conceptRef{IETF}{Internet Engineering Task Force (IETF)} es va crear el 1986
per desenvolupar i publicar estàndards per utilitzar a Internet. 
% 
El seu principal resultat són els documents \conceptRef{RFC}{``Request for Comment'' (RFC)},
que contenen descripcions tècniques completes d'aquests estàndards.


Aquests RFCs els creen grups de treball d'experts, 
perquè tothom els pugui utilitzar lliurement i de franc. 
% 
Qualsevol persona pot participar en aquests grups de treball, inclosos investigadors,
governs i agents de la indústria. 
% 
Les decisions es prenen quan hi està d'acord entre el 80 i el 90\% dels participants, de manera que cap 
entitat controla per si sola el procés d'estandardització.
% 
\begin{remark}
``Rebutgem els reis, els presidents i les votacions. 
Creiem en el consens aproximat i en el codi que funciona'' 
--- Dave Clark (IETF), 1992.

Consulta el DOI \href{https://dx.doi.org/10.1109/MAHC.2006.42}{\underline{10.1109/MAHC.2006.42}} per a 
una breu història de lluita pel poder.
\end{remark}

Els estàndards RFC són descriptius, no prescriptius. Això vol dir que ningú vigila l'ús correcte 
d'aquests estàndards. En canvi, els fabricants tenen l'incentiu de complir els RFCs 
perquè els seus productes siguin compatibles amb els d'altres. 
% 
El valor dels RFCs (i de qualsevol altre estàndard) depèn que
els usuaris els adoptin. Per aquest motiu, els RFCs es poden actualitzar i \conceptRef{retire (standard)}{retirar}.

Els protocols TCP/IP s'imposen perquè són flexibles i funcionen raonablement bé en la majoria de situacions, 
tot i que no són òptims pràcticament en cap. 
A més, els protocols TCP/IP són lliures, a diferència dels protocols \conceptRef{privative}{privatius}
que obligaven a comprar la solució d'un fabricant concret.
% 
\begin{remark}
Avui és difícil imaginar un món \textit{no} 
dominat per l'arquitectura TCP/IP, 
però es van comercialitzar i mantenir piles completament diferents fins a finals dels noranta.
En són exemples AppleTalk i IPX/SPX, d'Apple i Novell, respectivament.
\end{remark}

% https://www.icann.org/en/blogs/details/cheers-to-the-multistakeholder-community-30-9-2016-en

Alguns aspectes de TCP/IP, com ara els identificadors únics 
(inclòs el \concept{DNS}, el sistema de noms de domini)
i els números reservats, han de ser acordats per tothom perquè Internet funcioni correctament.
% 
La \concept{IANA} (Internet Assigned Numbers Authority) treballa en coordinació amb la \concept{IETF}
per definir els \conceptRef{RFC}{RFCs} adequats.

\begin{remark}
Des del 2016, la IANA és gestionada per una organització sense ànim de lucre amb múltiples parts interessades 
(\conceptRef{ICANN}{ICANN - Internet Corporation for Assigned Names and Numbers}).
% 
Des del 2004, la responsabilitat dels dominis regionals 
(\eg, \url{.cat}, \url{.fr}, \url{.es}) s'ha transferit a organitzacions 
d'àmbit continental anomenades registres regionals d'Internet (\conceptRef{RIR}{RIRs}).
\end{remark}



% # Based on https://www.youtube.com/watch?v=oIezCGjxV3A
% year,text
% 1962,Packet Switching
% 1969,ARPANET
% 1974,TCP/IP
% 1986,IETF
% 1991,WWW
% 1993,CIDR addressing
% 1998,Google
% ``The Internet doesn't do what the endpoints can do (\eg, reliability or compression).''


% - In February 1980, the Institute of Electrical and Electronics Engineers (IEEE) started project 802 to standardize local area networks (LAN).[16][27 (https://en.wikipedia.org/wiki/Ethernet)

% "protocols": "the first sheet of a volume" (on which contents and errata were written

% \begin{remark}
% Extra:
% \begin{itemize}
%     \item Video: https://www.youtube.com/watch?v=oIezCGjxV3A (part 1), https://www.youtube.com/watch?v=RN4gSBTANUY (part 2)
%     \item Monographic: https://cacm.acm.org/research/on-the-hourglass-model/
% \end{itemize}
% \end{remark}


% \begin{itemize}
% \item [\textbf{Layer 1}:] \textbf{Physical}
%     \begin{itemize}
%     \item[Features:] Ship \concept{digital} \otherBase{0}s and \otherBase{1}s through the real, \concept{analog} world.
%     \item[Offers:] Send and receive blocks of bytes.
%     \end{itemize}
% 
% \item [\textbf{Layer 2}:] \textbf{Network}\\
% 
% 
% \item [\textbf{Layer 3}:] \textbf{Internet}\\
% \item [\textbf{Layer 4}:] \textbf{Transport}\\
% \item [\textbf{Layer 7}:] \textbf{Application}\\
% \end{itemize}
% 
% \begin{remark}
% The \concept{OSI} model is a conceptual stack of layers that divides problems a bit differently from \concept{TCP/IP}. 
% TCP/IP was created on the go while the infrastructure of the Internet started its development, became its de-facto standard,
% and is used everywhere. The OSI model is useful to analyze and design communication protocols, but has limited practical usage.
% \end{remark}
% 
% 
% The network stack is normally part of the \concept{operating system} (OS).
% Only the \inlineCode{send} and \inlineCode{receive} functions of Layer 4 are exposed to the users of the OS.
% %
% This way, data transmitted by the \concept{network interface} are guaranteed to be properly formatted and follow
% the protocols of each stack layer.
% 
% \begin{remark}
% In order to control the lower layers --\eg, to transmit custom PDUs--, one needs to be the \texttt{root} user
% or give the \texttt{CAP\_NET\_RAW} kernel capability to a user program
% (or the equivalent in privative systems).
% \end{remark}
% 
% 
% \section{}\label{sec:stack_layers:organization}
% 
% - internet organization 
% bodies and stuff
